% Unidad U016. Registro firmado, lectura terminal y reconstruccion de K.
% Redaccion autonoma desde la autoridad material 1063 y su sucesor V2.

\chapter{Signed dodecaphasic register and integral reconstruction of the seal}
\label{chap:registro-dodecafasico-hadamard-k}

The construction of the integer seal does not begin with \(\alpha\), its inverse
or a metrological table. Its domain is the enriched terminal state
of the Topological Phase Kernel obtained after the already constructed chain
\begin{equation}
 \boxed{
 \mathrm{APP}\longrightarrow\mathrm{TRIT}\longrightarrow\mathrm{TPK}
 \longrightarrow\Omega_{\rm seed}
 \xrightarrow{T_{\min}}\mathcal U_6^{\rm cat}
 \xrightarrow{\rho_3}W_{243}
 \longrightarrow w_{30}\longrightarrow R_{36}
 \longrightarrow G_9\longrightarrow x_{\rm term}.}
 \label{eq:cadena-upstream-registro-dodecafasico}
\end{equation}
The foregoing arrows were defined with their respective domains:
pairs of seeds, signatures, emissions, ternary words, boundaries,
compatible histories, and nonadic holonomy. This unit does not reselect
any of those data. It constructs the integral charts of the same
terminal state and proves their equivalence.

\section{Windows of the oriented trace of length one hundred eight}

Let
\[
 \Theta_{108}=\mathbb Z/108\mathbb Z
\]
the observable calendar. Its twelve nonadic windows are
\[
 E_m=\{9m+1,\ldots,9m+9\},
 \qquad 0\leq m<12.
\]
The partition is disjoint and exhaustive:
\[
 \{1,\ldots,108\}=\bigsqcup_{m=0}^{11}E_m.
\]
The set of address codes that publish an event is fixed as
\[
 D_{\rm evt}=\{2,4,6,8\}.
\]
The orientation of the twelve windows is
\begin{equation}
 \Sigma_{12}=(+,+,+,-,-,-,+,+,+,-,-,-).
 \label{eq:signatura-doce-ventanas}
\end{equation}
If \(d_t\) is the direction encoded at step \(t\), the signed event
of window \(m\) is
\begin{equation}
 e_m=\Sigma_{12}(m)\,
 \#\{t\in E_m:d_t\in D_{\rm evt}\}.
 \label{eq:evento-firmado-ventana}
\end{equation}
The formula distinguishes the nonnegative number of events from the sign of the
orientation. It does not replace a cardinal direction with a sign: both
coordinates remain present in the TPK state.

\begin{definition}[Enhanced dodecaphonic registration]
\label{def:registro-enriquecido-dodecafasico}
The record transported by the trace is
\[
 \mathcal R_{12}(x)=
 \bigl((E_m,\tau_m,\sigma_m,\kappa_m,\Lambda_m)\bigr)_{m=0}^{11},
\]
where \(\tau_m\) is the TRIT orientation, \(\sigma_m\) the boundary signature, \(\kappa_m\) the carry, and \(\Lambda_m\) the ledger segment preserving the route. There is a chain of projections
\[
 \mathcal X_{\rm TPK}^{\rm enr}
 \longrightarrow\mathcal R_{12}
 \longrightarrow\Theta_{108},
\]
but neither of the two arrows is an identification.
\end{definition}

In particular, \(\Theta_{108}\) preserves only the position in the calendar.
The record \(\mathcal R_{12}\) additionally preserves orientation, signature,
carry, and provenance. The complete state further preserves cylinder,
boundary, quotient, terminal, and vertical memory.

\section{Integral chart of the dodecaphasic opposition}

For \(0\leq i<6\), the opposite positions determine
\begin{equation}
 p_i=e_i+e_{i+6},
 \qquad
 a_i=e_i-e_{i+6}.
 \label{eq:pares-opuestos-registro}
\end{equation}
The central charge and the five relative margins are
\begin{equation}
 s_C=\sum_{i=0}^{5}p_i,
 \qquad
 M_i=p_i-p_5\quad(0\leq i<5).
 \label{eq:carga-margenes-registro}
\end{equation}

\begin{definition}[Integral chart \(1+5+6\)]
\label{def:carta-integral-156}
It is defined
\begin{equation}
 \mathcal L_{12}(e_0,\ldots,e_{11})
 =\bigl(s_C,M_0,\ldots,M_4,a_0,\ldots,a_5\bigr).
 \label{eq:carta-integral-156}
\end{equation}
The coordinate partition is
\[
 1+5+6=12.
\]
\end{definition}

\begin{theorem}[Exact inverse of the integral chart]
\label{thm:inversion-carta-integral-156}
A tuple
\[
 \ell=(s_C,M_0,\ldots,M_4,a_0,\ldots,a_5)\in\mathbb Z^{12}
\]
belongs to the image of \(\mathcal L_{12}\) if and only if
\begin{align}
 s_C-\sum_{i=0}^{4}M_i&\equiv0\pmod6,
 \label{eq:congruencia-seis-carta}\\
 p_i&\equiv a_i\pmod2
 \quad(0\leq i<6),
 \label{eq:congruencia-paridad-carta}
\end{align}
where
\begin{align}
 p_5&=\frac{s_C-\sum_{i=0}^{4}M_i}{6},
 \label{eq:inversa-p5-carta}\\
 p_i&=p_5+M_i\quad(0\leq i<5).
 \label{eq:inversa-pi-carta}
\end{align}
In that case, the only entire preimage is
\begin{equation}
 e_i=\frac{p_i+a_i}{2},
 \qquad
 e_{i+6}=\frac{p_i-a_i}{2}.
 \label{eq:inversa-eventos-carta}
\end{equation}
\end{theorem}

\begin{proof}
Summing \(M_i=p_i-p_5\) for \(0\leq i<5\), we obtain
\[
 \sum_{i=0}^{4}M_i
 =\sum_{i=0}^{4}p_i-5p_5
 =s_C-6p_5,
\]
which forces \eqref{eq:inversa-p5-carta} and the congruence modulo six.
Once the six \(p_i\) have been reconstructed, the system
\[
 p_i=e_i+e_{i+6},
 \qquad a_i=e_i-e_{i+6}
\]
has the unique solution \eqref{eq:inversa-eventos-carta}. This solution is
integral exactly when \(p_i\) and \(a_i\) have the same parity.
The two families of congruences are therefore necessary and sufficient.
\end{proof}

This theorem precedes the numerical closure of \(\alpha\). It explains how the
oriented trace produces a reversible integral chart, and why preserving
the reduced calendar is not enough.

\section{Selection of the common terminal state}

The three compatible sections reach terminal depth with catalogs
of cardinalities
\[
 |\mathcal T_\pi|=196,
 \qquad
 |\mathcal T_e|=197,
 \qquad
 |\mathcal T_\varphi|=196.
\]
The product contains
\begin{equation}
 196\cdot197\cdot196=7\,567\,952
 \label{eq:cardinal-producto-terminal}
\end{equation}
ordered triples.

The column margin of the terminal reader is
\[
 q_\partial=(1,2,2,1,2,0)\in\mathbb F_3^6,
\]
and the row orientation of the combined channel is
\[
 \sigma=(1,1,-1)=(1,1,2)\in\mathbb F_3^3.
\]
The first datum acts on columns; the second acts on rows. Neither is deducible from the other.

The local signatures of Gram matrix, norm, weight, and distance reduce
\eqref{eq:cardinal-producto-terminal} to \(331\) triples. The margin
\(q_\partial\) leaves two indices:
\[
 (120,197,157),
 \qquad
 (129,197,148).
\]
Their matrices are
\[
 B_0=
 \begin{pmatrix}
 0&2&0&2&2&0\\
 0&0&1&2&2&0\\
 1&0&1&0&1&0
 \end{pmatrix},
 \qquad
 B_1=
 \begin{pmatrix}
 0&2&1&0&2&0\\
 0&0&1&2&2&0\\
 1&0&0&2&1&0
 \end{pmatrix}.
\]
The row sums are
\[
 (0,2,0),
 \qquad
 (2,2,1).
\]
Since
\[
 \langle\sigma\rangle
 =\{(0,0,0),(1,1,2),(2,2,1)\},
\]
only \(B_1\) satisfies the axial condition.

\begin{theorem}[Uniqueness of the Visible Terminal Publication]
\label{thm:unicidad-publicacion-terminal-visible}
The joint selection determines
\begin{equation}
 B_{\rm term}=
 \begin{pmatrix}
 021020\\001220\\100210
 \end{pmatrix},
 \label{eq:matriz-terminal-visible}
\end{equation}
corresponding to the triple \((129,197,148)\). No decimal evaluation of
\(\alpha\) enters the census or the two selection conditions.
\end{theorem}

\begin{proof}
The census and margin leave exactly \(B_0,B_1\). The charge of \(B_0\)
does not belong to \(\langle\sigma\rangle\), whereas that of \(B_1\) is
\(2\sigma\). The second candidate is therefore the only one that simultaneously satisfies
both typed requirements.
\end{proof}

\section{Two publications of the same state, not a conversion between them}

Let
\[
 x_{\rm term}\in
 \mathcal X_{\rm TPK}^{\rm term,enr}
 \subseteq\mathcal X_{\rm TPK}^{[108]}.
\]
The visible reading is
\[
 \pi_{\rm term}(x_{\rm term})=B_{\rm term}.
\]
The signed chart preserves another coordinate of the same state.

\begin{definition}[Signed reader produced by the terminal register]
\label{def:subespacio-firmado-terminal}
The reader's input state is exclusively the enriched terminal state already produced by the causal pipeline \(\mathcal U_t=\mathsf{Upd}_t\circ\mathsf{Tra}_t\circ\mathsf{Sel}_t\):
\begin{equation}
 x_{\rm term}
 =\bigl(\mathcal U_{108}\circ\cdots\circ\mathcal U_1\bigr)
   (x_{\rm seed})
 \in\mathcal X_{\rm TPK}^{\rm term,enr}.
 \label{eq:subespacio-firmado-terminal}
\end{equation}
Composing the oriented register with its channel readers \(90/120\) extracts from ledger \(\lambda_{108}\), without adding external coordinates,
\begin{equation}
 \mathcal E_{108}^{90,120}\!\left(\mathcal R_{12}(x_{\rm term})\right)
 =\bigl(b^{90},b^{120},Q_{\rm TPK}\bigr),
 \label{eq:canales-ledger-previos-u}
\end{equation}
where
\begin{align}
 b^{90}={}&(-495,-116,-684,108,601,-90,
             -173,-88,313,560,-397,461),\notag\\
 b^{120}={}&(-425,-281,-481,671,-255,30,
              475,-543,680,251,6,-128),\notag\\
 Q_{\rm TPK}={}&6263.\notag
\end{align}
These twenty-five coordinates are outputs of the register of the \cref{def:registro-enriquecido-dodecafasico}, not data added to the domain.

On a compatible output \((b^{90},b^{120},Q)\), we define the integral harmonic reader \(\Pi_H\) explicitly. For \(r=1,2,3\), let
\begin{equation}
 B_r=\sum_{j=0}^{3}b^{120}_{r+3j},
 \quad
 A_1=\frac{Q+2B_1+B_2}{3},
 \quad A_2=A_1-B_1,
 \quad A_3=A_2-B_2,
 \label{eq:lector-armonico-a-desde-ledger}
\end{equation}
and \(d_{r,j}=b^{90}_{r+3j}\). The orbital block \(r\) is
\begin{equation}
 \Pi_H^{(r)}(b^{90},b^{120},Q)
 =\bigl(A_r,d_{r,1}-d_{r,3},
              d_{r,0}+d_{r,2},d_{r,0}-d_{r,2}\bigr).
 \label{eq:lector-armonico-bloque-desde-ledger}
\end{equation}
Interleaving the three blocks by columns yields \(\Pi_H\), and the signed reading is constructed as the composition
\begin{equation}
 \boxed{
 R_{\rm sgn}:=
 \Pi_H\circ\mathcal E_{108}^{90,120}\circ\mathcal R_{12}:
 \mathcal X_{\rm TPK}^{\rm term,enr}
 \longrightarrow\mathcal U_{12}^{\rm sgn}:=
 \operatorname{im}R_{\rm sgn}\subseteq\mathbb Z^{12}.}
 \label{eq:lector-firmado-terminal}
\end{equation}
Neither \(U\) nor \(K\) belongs to the domain: \(U\) is the output of \eqref{eq:lector-firmado-terminal}, and \(K\) is subsequently reconstructed through the integral Hadamard inverse.
\end{definition}

The correct joint writing is
\begin{equation}
 x_{\rm term}
 \xrightarrow{(\pi_{\rm term},R_{\rm sgn})}
 (B_{\rm term},U_{12}^{\rm sgn}).
 \label{eq:doble-publicacion-terminal}
\end{equation}
No arrow \(B_{\rm term}\to U_{12}^{\rm sgn}\) is interposed: the visible matrix has forgotten precisely the orientation, opposition, quotient, carry, and ledger coordinates used by the signed chart. Both publications originate in the enriched state produced by \eqref{eq:cadena-upstream-registro-dodecafasico}.

The evaluation of the signed record gives
\begin{equation}
 \boxed{
 U=(2378,1406,2479,-452,998,-551,
 -668,-204,-371,-322,-28,-997).}
 \label{eq:vector-u-firmado-autonomo}
\end{equation}

\section{Naturality with respect to depth}

For \(N'\geq N\geq108\), the truncation
\begin{equation}
 \tau_{N',N}:
 \mathcal X_{\rm TPK}^{[N'],\rm enr}
 \longrightarrow\mathcal X_{\rm TPK}^{[N],\rm enr}
 \label{eq:truncamiento-preserva-ku-autonomo}
\end{equation}
acts on the enriched history and its ledger, not on a pair \((K,U)\) added to the state.

\begin{theorem}[Square of naturality of the signed reading]
\label{thm:naturalidad-lector-firmado-autonomo}
We have
\begin{equation}
 \boxed{
 R_{{\rm sgn},N}\circ\tau_{N',N}
 =R_{{\rm sgn},N'}.}
 \label{eq:cuadrado-naturalidad-lector-firmado}
\end{equation}
\end{theorem}

\begin{proof}
Each updater \(\mathcal U_t\) adds the next incidence to the ledger without rewriting the one hundred eight windows already generated. Therefore \(\mathcal R_{12}\), \(\mathcal E_{108}^{90,120}\), and \(\Pi_H\) read the same coordinates when truncating any subsequent prolongation. Commutativity is an equality of maps on produced states, not the tautological preservation of an output included in the input.
\end{proof}

\section{Step-three orbits and Hadamard transform}

The cyclic order of the seal is distributed over the three orbits.
\begin{align}
 K^{(1)}&=(K_1,K_4,K_7,K_{10}),
 \label{eq:orbita-k-1}\\
 K^{(2)}&=(K_2,K_5,K_8,K_{11}),
 \label{eq:orbita-k-2}\\
 K^{(3)}&=(K_3,K_6,K_9,K_{12}).
 \label{eq:orbita-k-3}
\end{align}
Let
\begin{equation}
 H_4=
 \begin{pmatrix}
 1&1&1&1\\
 1&1&-1&-1\\
 1&-1&1&-1\\
 1&-1&-1&1
 \end{pmatrix}.
 \label{eq:matriz-hadamard-cuatro-autonoma}
\end{equation}
If \(P_{\rm orb}\) reorders the twelve coordinates according to
\eqref{eq:orbita-k-1}--\eqref{eq:orbita-k-3}, the global transformation is
\begin{equation}
 \mathcal H_{12}
 =P_{\rm orb}^{-1}
 (H_4\oplus H_4\oplus H_4)P_{\rm orb}.
 \label{eq:hadamard-global-doce}
\end{equation}

The three blocks of \eqref{eq:vector-u-firmado-autonomo}, read by
orbits, are
\begin{align*}
 U^{(1)}&=(2378,-452,-668,-322),\\
 U^{(2)}&=(1406,998,-204,-28),\\
 U^{(3)}&=(2479,-551,-371,-997).
\end{align*}

\begin{lemma}[Algebraic Hadamard quarter-turn]
\label{lem:cuarto-inversa-hadamard}
The matrix \(H_4\) satisfies
\begin{equation}
 H_4^2=4I_4,
 \qquad
 H_4^{-1}=\frac14H_4,
 \qquad
 \det H_4=-16.
 \label{eq:identidades-hadamard-cuatro}
\end{equation}
\end{lemma}

\begin{proof}
The rows of \(H_4\) are orthogonal, and each has squared norm four.
Since the matrix is symmetric, \(H_4^{\mathsf T}H_4=H_4^2=4I_4\). The
formula for the inverse and the determinant follow immediately.
\end{proof}

\begin{theorem}[Unique full reconstruction of the seal]
\label{thm:reconstruccion-entera-k-hadamard}
The inverse
\begin{equation}
 K^{(r)}=\frac14H_4U^{(r)}
 \qquad(r=1,2,3)
 \label{eq:inversion-bloques-hadamard}
\end{equation}
yields
\begin{align*}
 K^{(1)}&=(234,729,621,794),\\
 K^{(2)}&=(543,659,058,146),\\
 K^{(3)}&=(140,824,914,601).
\end{align*}
By undoing the orbital order one obtains
\begin{equation}
 \boxed{
 K=(234,543,140,729,659,824,621,058,914,794,146,601).}
 \label{eq:sello-k-doce-autonomo}
\end{equation}
It is the unique rational preimage of \(U\), and it belongs to
\(\mathbb Z^{12}\).
\end{theorem}

\begin{proof}
By \eqref{eq:identidades-hadamard-cuatro}, the rational preimage is unique. The explicit multiplications are
\begin{align*}
 H_4U^{(1)}&=(936,2916,2484,3176),\\
 H_4U^{(2)}&=(2172,2636,232,584),\\
 H_4U^{(3)}&=(560,3296,3656,2404).
\end{align*}
All coordinates are divisible by four. Dividing and reinterleaving
yields \eqref{eq:sello-k-doce-autonomo}.
\end{proof}

The factor \(1/4\) is not chosen to fit an output: it is the inverse
forced by \(H_4^2=4I_4\). Nor is it the square root of \(-1\). The
internal complex structure \(A_W^2=-I\) will appear after constructing the
Paley matrix; these are two facts of different types.

\section{Integral structure and Smith normal form}

\begin{theorem}[Integral image of \(H_4\)]
\label{thm:smith-hadamard-autonomo}
The Smith normal form is
\begin{equation}
 \operatorname{SNF}(H_4)=\operatorname{diag}(1,2,2,4).
 \label{eq:snf-hadamard-autonoma}
\end{equation}
Consequently,
\begin{equation}
 \operatorname{coker}\mathcal H_{12}
 \cong(\mathbb Z/2\mathbb Z)^6
 \oplus(\mathbb Z/4\mathbb Z)^3,
 \label{eq:cociente-hadamard-doce}
\end{equation}
and the image of \(\mathcal H_{12}\) has index
\[
 16^3=4096
\]
in \(\mathbb Z^{12}\). For \(u\in\mathbb Z^4\),
\begin{equation}
 u\in H_4\mathbb Z^4
 \quad\Longleftrightarrow\quad
 H_4u\equiv0\pmod4.
 \label{eq:criterio-imagen-hadamard}
\end{equation}
\end{theorem}

\begin{proof}
The greatest common divisor of the entries is one. The minors of order two are
even, and one has absolute value two; those of order three are multiples of
four, and one has absolute value four; the determinant has modulus
sixteen. The determinantal divisors are \(1,2,4,16\), yielding
the invariant factors \(1,2,2,4\). The triple direct sum repeats the
factors and multiplies the indices. Finally,
\(H_4^{-1}=H_4/4\), so the preimage is integral exactly under the
congruence \eqref{eq:criterio-imagen-hadamard}.
\end{proof}

\section{Certification in the second coordinate system: step-three and step-four differences}

With cyclic indices modulo twelve, we define
\begin{equation}
 (D_3K)_m=K_m-K_{m+3},
 \qquad
 (D_4K)_m=K_m-K_{m+4},
 \qquad
 Q(K)=\sum_{m=1}^{12}K_m.
 \label{eq:extractor-diferencias-k}
\end{equation}
Applied to the seal already reconstructed by Hadamard, these operators must recover exactly the channel coordinates produced by the TPK register before \(U\). For \eqref{eq:sello-k-doce-autonomo},
\begin{align*}
 D_3K={}&(-495,-116,-684,108,601,-90,\\
         &\qquad-173,-88,313,560,-397,461),\\
 D_4K={}&(-425,-281,-481,671,-255,30,\\
         &\qquad475,-543,680,251,6,-128),\\
 Q(K)={}&6263.
\end{align*}
Comparison with \eqref{eq:canales-ledger-previos-u} proves the equality of outputs
\begin{equation}
 \boxed{
 D_3K=b^{90},\qquad D_4K=b^{120},\qquad Q(K)=Q_{\rm TPK}.}
 \label{eq:certificacion-canales-ledger-desde-k}
\end{equation}
The differential system certifies and reconstructs the previous register; it does not select \(K\), \(U\), or any constant.

\begin{theorem}[Completeness of the transversal extractor]
\label{thm:completitud-extractor-transversal-k}
The operator
\[
 \mathcal E_{\rm dod}=(D_3,D_4,Q):
 \mathbb Z^{12}\longrightarrow\mathbb Z^{25}
\]
has rank twelve and Smith normal form
\begin{equation}
 \operatorname{SNF}(\mathcal E_{\rm dod})
 =\operatorname{diag}(\underbrace{1,\ldots,1}_{11},12).
 \label{eq:snf-extractor-transversal-k}
\end{equation}
In particular, \((D_3K,D_4K,Q(K))\) determina a single \(K\).
\end{theorem}

\begin{proof}
If \(D_3v=D_4v=0\), then \(v\) is constant along steps
three and four. Since \(\gcd(3,4)=1\), those steps generate all of
\(\mathbb Z/12\mathbb Z\), so that
\[
 \ker D_3\cap\ker D_4=\langle\mathbf1\rangle.
\]
The total charge does not vanish on that line, since \(Q(\mathbf1)=12\). The rank
is twelve. The rows of \((D_3,D_4)\) are oriented incidences of a connected
Cayley graph; a spanning tree supplies eleven unit
invariant factors. Every maximal minor must contain the row \(Q\), and the
unimodular operations reduce its final contribution to twelve. There is a maximal minor of
modulus twelve; therefore, the last factor is exactly twelve.
\end{proof}

The reconstruction of \(U\) from these coordinates does not need to
multiply again by \(H_4\). For \(r=1,2,3\), let
\begin{equation}
 B_r=\sum_{j=0}^{3}(D_4K)_{r+3j}.
 \label{eq:br-diferencias-k}
\end{equation}
Then
\[
 (B_1,B_2,B_3)=(972,-1073,101).
\]
Orbital sums are
\begin{align}
 A_1&=\frac{Q+2B_1+B_2}{3}=2378,
 \label{eq:a1-reconstruccion-u}\\
 A_2&=A_1-B_1=1406,
 \label{eq:a2-reconstruccion-u}\\
 A_3&=A_2-B_2=2479.
 \label{eq:a3-reconstruccion-u}
\end{align}
If \(d_{r,j}=(D_3K)_{r+3j}\), the signed block of orbit \(r\) is
\begin{equation}
 \bigl(A_r,
 d_{r,1}-d_{r,3},
 d_{r,0}+d_{r,2},
 d_{r,0}-d_{r,2}\bigr),
 \label{eq:bloque-firmado-desde-diferencias}
\end{equation}
that reproduces the three blocks \(U^{(r)}\).

\section{Threshold Tetrad of the Seal}

The local decimal completion satisfies
\[
 1000=729+271.
\]
The support of the seal coordinates that reach the threshold \(729\) is
\begin{equation}
 \boxed{
 B_K:=\{m:K_m\geq729\}=\{4,6,9,10\}.}
 \label{eq:tetrada-umbral-k}
\end{equation}
This tetrad is obtained after constructing \(K\). In this unit it is not yet
called a Witt star: that object will be defined only after
constructing the ternary code, its hexads, and the Steiner system.

\section{Causal independence with respect to the physical output}

\begin{theorem}[Isolation of \(U\) and \(K\)]
\label{thm:aislamiento-u-k-alpha-codata}
The composition
\begin{equation}
 \Omega_{\rm seed}\longrightarrow x_{\rm term}
 \xrightarrow{\mathcal R_{12}}
 \mathcal R_{12}(x_{\rm term})
 \xrightarrow{\mathcal E_{108}^{90,120}}
 (b^{90},b^{120},Q_{\rm TPK})
 \xrightarrow{\Pi_H}U
 \xrightarrow{\mathcal H_{12}^{-1}}K
 \label{eq:composicion-upstream-u-k}
\end{equation}
does not receive \(\alpha\), \(\alpha^{-1}\), CODATA or a target decimal
string as input. The seal \(K\) is determined before defining the
recurrence that will produce \(\alpha\).
\end{theorem}

\begin{proof}
The first part of the composition uses only seeds, emissions, ternary readers, operators \(L_0,L_1\), survival relations, phase transport, and TPK state charts. Register \(\mathcal R_{12}\) and its channel reader produce \((b^{90},b^{120},Q_{\rm TPK})\); formulas \eqref{eq:lector-armonico-a-desde-ledger} and \eqref{eq:lector-armonico-bloque-desde-ledger} produce \(U\); only then does the linear transformation \(\mathcal H_{12}^{-1}=\mathcal H_{12}/4\) act on its integral image. None of the formulas contains a physical variable or a metrological table. The base-thousand recurrence is introduced only in the next unit, after \(K\) has been fixed.
\end{proof}

\section{Refutation criteria}

\begin{criterion}[Refutation of the register and the seal]
The construction is refuted if any of the following facts
is established:
\begin{enumerate}
 \item the windows \(E_m\) do not partition the one hundred eight steps;
 \item the chart \(\mathcal L_{12}\) fails its inverse or its congruences;
 \item the terminal census does not produce \(331\) triples, two matrices and a unique
       axial charge;
 \item an attempt is made to obtain \(U\) from \(B_{\rm term}\) after forgetting the
       signed coordinates;
 \item the naturality square
       \eqref{eq:cuadrado-naturalidad-lector-firmado} fails;
 \item \(H_4^2\neq4I_4\), some preimage is not integral or the
       reconstructed vector differs from \eqref{eq:sello-k-doce-autonomo};
 \item the Smith normal form differs from
       \eqref{eq:snf-hadamard-autonoma};
 \item extractor \((D_3,D_4,Q)\) loses rank or fails to reproduce \(U\);
 \item \(B_K\) is introduced before constructing \(K\), or is assigned a
       Witt structure before defining \(W_{12}\);
 \item a digit of \(\alpha\) or a metrological reference intervenes in
       any of the arrows of \eqref{eq:composicion-upstream-u-k}.
\end{enumerate}
\end{criterion}
